%% ============================================================================
%% AISec 2026 submission — assembled from five section drafts.
%% ASSEMBLED, NOT AUTHORED: prose is integrated verbatim from the drafts with
%% only number/cross-reference/anonymization edits. See AISEC-INTEGRATION-SUMMARY.md
%% for the full reconciliation table and the list of open author decisions.
%%
%% ANONYMIZATION (AISec 2026 is DOUBLE-BLIND — "properly anonymized ... no way to
%% identify authors, including when linking code repositories"):
%%   - review/anonymous class options are ON.
%%   - All GitHub URLs / author names / commit hashes are removed.
%%   - The tool name is behind \sysname. If the public repo makes the real name
%%     de-anonymizing, redefine \sysname to a neutral placeholder before submission.
%%
%% BUILD NOTE: requires the ACM `acmart` class (https://www.acm.org/publications/
%% proceedings-template). Not vendored here (no network); place acmart.cls beside
%% this file before compiling.
%% ============================================================================
\documentclass[sigconf,anonymous,review]{acmart}

\usepackage{booktabs}

%% ---- Anonymization-safe system name. Flip to a neutral name if 'PromptSonar'
%% is discoverable and de-anonymizing (see summary doc, open decision A2). ----
\newcommand{\sysname}{PromptSonar}

\settopmatter{printacmref=false}
\setcopyright{none}
\acmConference[AISec '26]{19th ACM Workshop on Artificial Intelligence and Security}{2026}{}
\acmYear{2026}
\acmDOI{}
\acmISBN{}

\begin{document}

\title{Reference-Verified Cross-File Execution-Path Analysis for AI-Agent Repositories}

%% Double-blind: authors anonymized for review.
\author{Anonymous Author(s)}
\affiliation{%
  \institution{Submission under double-blind review}
  \country{}
}

\begin{abstract}
Static analysis tools for AI-enabled software increasingly target individual artifacts --- a single prompt, a single agent ``skill,'' a single Model Context Protocol (MCP) server configuration --- in isolation. This artifact-level granularity is structurally unable to detect a class of risk that only exists across artifact boundaries: whether an instruction in one file can actually reach a privileged capability declared in another. We present \sysname{}, a deterministic, repository-scale static analyzer that constructs a reference-verified execution-path graph across prompts, agent skills, MCP configurations, workflows, and tool definitions, and evaluate it against 22 controlled benchmark cases and a 37-repository corpus of real-world AI-agent codebases. Compared against a representative artifact-level scanner on an identical corpus, \sysname{} identifies hundreds of cross-file execution paths that per-artifact analysis cannot observe by construction, including risk in 17 repositories where the baseline surfaces no connective finding at all. We additionally report a rigorous, audit-driven false-positive analysis: an initial 85\% false-positive rate (verified against source on a stratified, seeded sample) was reduced across three targeted root-cause fixes to 60\%, with true positives increasing from 3 to 8 in the same sample, and we report the full remaining false-positive taxonomy rather than an unaudited prevalence claim. We argue this audit methodology, and the honest disclosure of a still-substantial residual error rate, is itself a needed contribution for a nascent tool category where unverified precision and recall claims are otherwise the norm.
\end{abstract}

\maketitle

\section{Introduction}

The rapid adoption of autonomous AI agents --- systems that combine a natural-language prompt, a set of declared tools, and external capability providers via protocols such as MCP --- has created a new class of software artifact with no direct analogue in traditional application security. A prompt is not source code in the sense a conventional static analyzer parses; an MCP server configuration is not a dependency manifest in the sense a supply-chain scanner audits; yet together, a prompt, a skill definition, and an MCP configuration can compose into a system where a user- or model-generated instruction reaches shell execution, filesystem access, or credential exfiltration, without any single one of those three artifacts containing an obviously dangerous pattern in isolation.

Existing static analysis tooling for this space has emerged rapidly over the past year but has converged, largely, on artifact-level granularity: scanning a single skill definition, a single MCP configuration, or a single prompt for known-risky patterns (Section~\ref{sec:related}). This granularity is adequate for detecting a capability \emph{declared} in one place, but structurally cannot detect a capability \emph{reached} across files --- because doing so requires constructing a reference graph across the repository, which single-artifact tools do not build.

We present \sysname{}, a deterministic static analyzer that treats the AI-relevant portion of a repository --- prompts, agent skill definitions, MCP server configurations, workflow definitions, memory configurations, and tool implementations --- as a connected artifact graph, and identifies whether a reachable path exists from an instruction to a sensitive action, gated on verified references (imports, named configuration keys, explicit links) rather than incidental textual co-location. \sysname{} performs no LLM calls during analysis: all detection is pattern- and reference-based, which we verify produces byte-identical output across repeated runs (Section~\ref{sec:determinism}).

In developing and evaluating \sysname{}, we found that our own early cross-file reachability claims --- computed via a looser, capability-word-based edge-construction heuristic --- were substantially unreliable under manual audit (an initial 85\% false-positive rate on a stratified, source-verified sample). Rather than report the more impressive-looking unaudited numbers, we treat this finding as a first-class result: we describe the audit methodology, the root-cause-driven fix process that reduced (but did not eliminate) this false-positive rate, and the taxonomy of errors that remain, arguing that this kind of disclosure is currently rare in this tool category and is itself a contribution reviewers and practitioners should expect going forward.

\paragraph{Contributions.}
\begin{enumerate}
\item A reference-gated, cross-file execution-path analysis technique for AI-agent repositories that identifies reachability across prompt, skill, MCP, workflow, and tool boundaries --- a detection capability we show, via head-to-head evaluation against a representative artifact-level scanner on an identical 37-repository corpus, is structurally unavailable to single-artifact analysis regardless of within-file detection quality (Section~\ref{sec:crossfile}).
\item A deterministic, LLM-free implementation verified to produce reproducible output across repeated runs (Section~\ref{sec:determinism}), enabling secure analysis of sensitive repositories without an external inference dependency.
\item A rigorous false-positive audit methodology --- stratified, seeded, source-verified sampling with per-repository clustering controls --- applied across three development iterations, reducing false-positive rate from 85\% to 60\% via targeted, regression-tested root-cause fixes, with the complete remaining error taxonomy reported rather than suppressed (Section~\ref{sec:audit}, Section~\ref{sec:precision}).
\item A controlled benchmark of 22 cases with ground-truth annotations, released as part of our evaluation artifact, covering cross-file reachability, contextual severity adjustment, and true-negative behavior.\footnote{We are not aware of an equivalent public benchmark at this analysis granularity; this novelty claim should be re-verified against adjacent work before final submission (open decision, see summary).}
\end{enumerate}

\section{Related Work}
\label{sec:related}

We organize related work into three categories by analysis granularity, since this dimension --- not detection technique --- is what most directly determines whether a tool can observe cross-file execution reachability at all.

\subsection{Prompt- and content-level scanners}
A first generation of tools operates on individual prompts or model inputs/outputs in isolation: pattern-matching or classifier-based detection of prompt injection, jailbreak attempts, and PII/secret leakage within a single string. Tools in this category (e.g., garak~\cite{garak}, LLM Guard~\cite{llmguard}, NeMo Guardrails~\cite{nemoguardrails}, Rebuff~\cite{rebuff}) are typically deployed as an inline runtime filter or a red-teaming harness rather than a static repository scanner. They have no concept of a ``repository'' as an analysis unit and consequently no mechanism for reasoning about whether a given prompt's instructions are reachable to a sensitive action defined in a separate file. \sysname{} differs in both deployment model (static, pre-deployment, repository-scoped) and analysis unit (the artifact graph, not the single string).

\subsection{Artifact- and skill-level static scanners}
A second, more recent generation performs static analysis over individual AI-relevant artifacts --- most commonly agent ``skill'' definitions (e.g., \texttt{SKILL.md} files), individual MCP server configurations, or single tool definitions. NVIDIA SkillSpector~\cite{skillspector} is representative of this category: a two-stage static analyzer (pattern/AST matching, with an optional LLM-based semantic evaluation stage) that scores individual skills against a taxonomy of vulnerability categories. Similarly, MCP-specific scanners such as \texttt{mcp-scan}~\cite{mcpscan} evaluate a server configuration's declared permissions and tool exposure; notably, \texttt{mcp-scan} additionally reasons about \emph{cross-server} tool-shadowing and ``cross-origin'' escalation among multiple connected MCP servers at runtime --- a form of cross-artifact reasoning, though one operating on live server tool descriptions rather than a static repository reference graph.

These tools share a structural limitation directly relevant to our contribution: with the partial exception of \texttt{mcp-scan}'s runtime cross-server checks, they evaluate each artifact independently and have no \emph{static} execution-path model connecting artifacts across a repository. This follows from the scanning unit itself being a single file, a single declared skill, or a set of live server descriptions, rather than a graph of resolved references between repository artifacts. In our head-to-head evaluation on a shared 37-repository wild corpus (Section~\ref{sec:crossfile}), the baseline reported zero cross-file execution paths on every repository, by construction, while \sysname{} identified reference-verified cross-file paths in the majority of repositories scanned --- including 17 repositories where the per-artifact baseline surfaced no connective risk at all between a prompt's stated intent and a tool's actual capability.

Broader application-security platforms extending into the AI-agent space (e.g., MCP-aware modules within general-purpose developer security tools, and IDE-integrated agent scanners recently introduced by major vendors) largely follow the same artifact-level pattern: strong within-file detection, no static cross-artifact execution model.

\subsection{Traditional static application security testing (SAST)}
General-purpose SAST tools (e.g., Semgrep~\cite{semgrep}) perform taint-flow analysis --- including cross-file taint tracking --- over conventional source code, targeting vulnerability classes like SQL injection, command injection, and XSS in application logic written by human developers. These tools are mature and, notably, already treat ``exclude tests, fixtures, and documentation from production analysis'' as a first-class rule primitive (Section~\ref{sec:audit} discusses why \sysname{}'s comparable exclusion logic required several iterations to reach a similar standard).

However, these tools have no model of AI-specific artifacts: a \texttt{SKILL.md} file, an MCP server configuration, or an agent system prompt is not source code in the sense a SAST taint tracker is designed to parse, and the ``sink'' of interest --- an instruction reaching a privileged capability exposed to an autonomous agent --- is a different threat model than a data value reaching an unsanitized database call. \sysname{}'s contribution is complementary rather than competing with this category: it extends taint-style reachability reasoning to the AI-artifact graph specifically (prompts, skills, MCP configurations, workflows, memory definitions) rather than to conventional application code.

\subsection{Positioning}
\label{sec:positioning}
\sysname{}'s distinguishing claim, evaluated empirically in Section~\ref{sec:results}, is not that it detects a wider variety of vulnerability patterns than existing tools --- several of the artifact-level and content-level tools above have mature, extensively-tuned pattern libraries developed over longer periods. The claim is narrower and, we argue, more structurally significant: reference-verified execution-path reachability \emph{across} the artifact boundaries that define a modern AI-agent system (prompt $\rightarrow$ skill $\rightarrow$ MCP server $\rightarrow$ tool $\rightarrow$ sensitive action) is a detection capability that artifact-isolated and content-isolated tools cannot provide regardless of how their within-file detection improves, because it requires an analysis unit --- the whole-repository reference graph --- that those tools do not construct.

\section{System Design}
\label{sec:design}
%% ======================================================================
%% GAP: The five provided drafts are Abstract+Intro (1), Related Work (2),
%% Methods (4), Results (5), Limitations (6). NO Section 3 (the technique /
%% system design) was provided, yet every cross-reference in the drafts
%% assumes Methods=4/Results=5/Limitations=6, i.e. that Section 3 exists.
%% This placeholder preserves the numbering. The author must supply the
%% System Design section (artifact classification, reference-graph
%% construction, reachability/edge gating, contextual verdict engine),
%% likely adaptable from docs/papers/mcp-security-arxiv-draft.md.
%% See summary doc, open decision D1 (highest-priority content gap).
%% ======================================================================
\emph{[Section~3 (System Design) is not among the supplied drafts and must be authored: artifact discovery/classification, the reference-verified edge-construction rule, reachability computation, and the contextual verdict/severity engine. Placed here to preserve the cross-reference numbering the other sections assume.]}

\section{Methods}
\label{sec:methods}

\subsection{System under test}
\label{sec:sut}
We evaluate \sysname{} at a fixed, pinned commit\footnote{Exact commit hash and artifact URL omitted for double-blind review; recorded in the (anonymized) evaluation artifact.} of a deterministic, repository-scale static analyzer for AI-enabled codebases. \sysname{} performs no LLM calls during scanning; all analysis is pattern- and reference-based, and we verify determinism directly (Section~\ref{sec:determinism}).

\subsection{Controlled corpus}
\label{sec:controlled-corpus}
The controlled corpus consists of 22 manually constructed test cases with ground-truth annotations covering: single-file findings with a directly-evidenced sensitive action; cross-file execution paths requiring a genuine reference between a prompt/skill and a tool or MCP configuration; MCP configurations under both permissive (\texttt{auto\_approve: true}) and restrictive control states; and true-negative cases (clean prompts and configurations expected to produce zero findings), included specifically to measure false-positive behavior on benign input as well as detection recall on planted issues.

\subsection{Wild corpus}
\label{sec:wild-corpus}
The wild corpus consists of 37 real-world, actively-maintained public repositories, shallow-cloned (depth 1) at evaluation time, spanning three categories: MCP server implementations, agent orchestration frameworks (including LangChain, AutoGen, CrewAI, pydantic-ai, and Semantic Kernel), and AI coding-assistant tooling (including Continue, Cline, and OpenHands). Repository selection prioritized popular, actively-maintained projects representative of production AI-agent development, described further as a limitation in Section~\ref{sec:eval-limits}. Cloned repository contents are excluded from version control; the repository list and clone script are included in our evaluation artifact for reproducibility, though we note that repository content may drift between clone events (Section~\ref{sec:eval-limits}), and camera-ready reproduction should pin exact commit SHAs per repository rather than relying on default-branch snapshots as we did during development iteration.\footnote{One repository (pydantic-ai) exposed an $O(n^2)$ path-construction regression during development that caused a timeout; the fix is included in the evaluated commit. Re-running on an earlier commit may not reproduce the reported 0 scan failures.}

\subsection{Determinism verification}
\label{sec:determinism}
We verified engine determinism by running identical scans three times against an identical repository state and diffing output after stripping run-scoped metadata fields (a generated report identifier and two timestamp fields, which are expected to vary between runs and are not analysis output). All substantive fields --- findings, severities, execution paths, and confidence scores --- were byte-identical across all three runs.

\subsection{Baseline: an artifact-level scanner}
\label{sec:baseline}
We compare against NVIDIA SkillSpector~\cite{skillspector}, a static analyzer for individual AI agent skill definitions, as the most directly comparable existing tool at the artifact-analysis granularity (Section~\ref{sec:related}). We run it in its deterministic mode (\texttt{--no-llm}) throughout, to keep the comparison free of LLM-induced non-determinism and API cost, and note where this excludes its optional LLM-assisted semantic evaluation stage from our comparison.

On the controlled corpus, we score the baseline against the subset of cases within its stated scope (single-artifact analysis) and report cases requiring cross-file or MCP-configuration context as out-of-scope rather than scoring them as failures, since it is not designed to evaluate them.

On the wild corpus, an initial whole-repository comparison caused the baseline to exceed our per-repository timeout on 18 of 37 repositories; we determined via direct timing measurement that this was caused by its per-dependency vulnerability (OSV) lookup stage on large repositories, not by its skill-analysis logic. We therefore scoped the baseline's wild-corpus comparison to each repository's AI-artifact surface (the same prompt/skill/agent/MCP files \sysname{} analyzes) rather than the full repository tree, which resolved all timeouts and allows a complete 37/37 comparison on the capability under test.

\subsection{False-positive audit protocol}
\label{sec:audit-protocol}
We audit \sysname{}'s wild-corpus output against source ground truth using the following fixed protocol, applied identically across all rounds reported in Section~\ref{sec:audit}:
\begin{enumerate}
\item \textbf{Sampling.} From the full population of wild-corpus findings (security issues) and reachable execution paths, draw a stratified random sample of 20 items using a fixed seed (20260703), targeting an approximate 3:10:7 split across Confirmed/Probable/Potential confidence tiers. A per-repository cap of 2 items is enforced across the full sample to prevent a single repository's characteristics from dominating the tier-level false-positive estimate.
\item \textbf{Verification.} For each sampled item, retrieve the actual source file(s) from the cloned repository and manually verify the claim against real content: for a security finding, whether the flagged text is genuinely security-relevant production content (as opposed to a comment, docstring, test fixture, or documentation example); for an execution path, whether each edge in the claimed path corresponds to a real reference (import, config key, named link) in the source, as opposed to an inferred connection based on incidental co-location of related terminology.
\item \textbf{Verdict.} Each item is labeled Confirmed, False Positive, Overstated (a real underlying issue exists but is mischaracterized), or Unable to Verify.
\item \textbf{Reporting.} We report the false-positive rate overall and by confidence tier, and categorize all false positives by root cause before deciding whether and how to address them, rather than fixing individual instances ad hoc.
\end{enumerate}
This protocol was applied across four audit rounds during development (Section~\ref{sec:audit}), each following a targeted fix addressing the previous round's dominant root cause, verified via regression tests to not reintroduce previously-fixed false-positive classes before re-auditing.

\subsection{Reproducibility}
\label{sec:repro}
Tool and dependency versions are pinned and recorded in the evaluation artifact, including the Node.js runtime, the \sysname{} commit, the baseline commit/version, and any auxiliary tooling. All evaluation scripts (corpus construction, scanning harness, table generation) are included. As noted in Section~\ref{sec:wild-corpus}, exact wild-corpus repository commit SHAs were not pinned during development iteration (default-branch snapshots were used); we pin exact SHAs for the camera-ready artifact to guarantee bit-for-bit reproducibility of the reported wild-corpus numbers.

\section{Results}
\label{sec:results}

\subsection{Controlled Evaluation}
\label{sec:controlled-eval}
We evaluated \sysname{} against a manually-constructed benchmark of 22 cases spanning single-file and cross-file execution paths, MCP configurations with varying approval states, and clean (true-negative) fixtures. Table~\ref{tab:controlled} reports precision, recall, and F1 against ground-truth annotations. Some cases carry multiple expected findings, so the true-positive count (25) exceeds the case count (22).

\begin{table}[t]\centering\caption{Controlled corpus: detection accuracy on 22 ground-truthed cases.}\label{tab:controlled}
\begin{tabular}{lrrrrrrr}\toprule
Tool & Cases & TP & FP & FN & Prec. & Rec. & F1 \\\midrule
\sysname{} & 22 & 25 & 0 & 0 & 1.00 & 1.00 & 1.00 \\
Baseline ($\neg$LLM) & 12$^{a}$ & 4 & 0 & 5 & 1.00 & 0.44 & 0.62 \\
Baseline (+LLM) & -- & -- & -- & -- & -- & -- & -- \\
\bottomrule\end{tabular}
\\[2pt]\footnotesize $^{a}$The baseline is skill/file-level; 10/22 cases (cross-file, MCP-config, memory) are out of its scope and excluded. (+LLM) not run: requires an API key and would violate the deterministic, no-cost constraint.
\end{table}


\sysname{} matched all 15 expected execution paths and correctly classified all 4 true-negative (clean) fixtures. We report perfect scores on this benchmark with an important caveat, addressed in Section~\ref{sec:eval-limits}: the controlled corpus was constructed by the same team that built the tool, and construct validity should not be assumed to transfer to the wild-corpus setting evaluated next.

\subsection{Wild-Corpus Evaluation}
\label{sec:wild-eval}
We ran \sysname{} against 37 real-world, actively-maintained repositories spanning MCP server implementations, agent frameworks, and coding-assistant tooling. Repositories were selected to represent the population of publicly available AI-agent codebases, not to maximize findings. The final wild-corpus scan (at the audited engine state, Section~\ref{sec:audit}) is summarized in Table~\ref{tab:wild}.

\begin{table}[t]\centering\caption{Wild corpus: aggregate over 37 public AI repositories (audited engine state).}\label{tab:wild}
\begin{tabular}{lr}\toprule Metric & Value \\\midrule
Repositories scanned & 37 / 37 \\
Total findings & 723 \\
\quad critical / high / medium / low & 0 / 30 / 15 / 678 \\
\quad confirmed / probable / potential & 25 / 512 / 186 \\
Execution paths & 1414 \\
Cross-file paths & 520 \\
Mean / median / max scan time (s) & 10.9 / 6.6 / 38.7 \\
\bottomrule\end{tabular}\end{table}


\subsection{Comparison to the Baseline: The Cross-File Gap}
\label{sec:crossfile}
We compared \sysname{} against the artifact-level baseline on the same 37-repository wild corpus, scoped fairly to each repository's AI-artifact surface (the same files \sysname{} analyzes) to keep whole-repository file-walk cost from timing the baseline out on large monorepos.

\begin{table}[t]\centering\caption{Cross-file detection gap. \sysname{} reconstructs reference-verified cross-file execution paths; the per-file baseline has no cross-file execution-path model and detects zero on every repository. Top 10 repositories by cross-file path count shown.}\label{tab:crossfile}
\begin{tabular}{lrr}\toprule
Repository & \sysname{} & Baseline \\\midrule
\texttt{langchain-ai/langchainjs} & 97 & 0 \\
\texttt{upstash/context7} & 94 & 0 \\
\texttt{cline/cline} & 90 & 0 \\
\texttt{RooCodeInc/Roo-Code} & 77 & 0 \\
\texttt{block/goose} & 50 & 0 \\
\texttt{anthropics/anthropic-cookbook} & 23 & 0 \\
\texttt{pydantic/pydantic-ai} & 18 & 0 \\
\texttt{yamadashy/repomix} & 18 & 0 \\
\texttt{Aider-AI/aider} & 12 & 0 \\
\texttt{openai/openai-agents-python} & 10 & 0 \\
\midrule
\textbf{Total (37 repos)} & \textbf{520} & \textbf{0} \\
\textbf{Repos with $\geq$1 cross-file path} & \textbf{17 / 37} & \textbf{0} \\
\bottomrule\end{tabular}
\\[2pt]\footnotesize The baseline scans per-file/per-skill and has no cross-file execution-path model; all \sysname{} cross-file paths are out of its detection scope. Compared on 37/37 repos (0 timeouts under AI-artifact scoping).
\end{table}


The baseline, by design, evaluates artifacts independently and has no static cross-file execution model --- it reported 0 cross-file paths on every repository in the corpus. \sysname{} identified 520 reachable execution paths spanning prompt/skill/MCP/tool boundaries, with at least one such path in 17 repositories (45.9\% of the corpus) where the baseline's per-file analysis found nothing connecting those risks (Table~\ref{tab:crossfile}). This is the paper's central empirical claim: instruction-to-sensitive-action reachability that crosses file boundaries is a real and currently under-served detection surface, and single-artifact scanners cannot see it by construction, not merely by current implementation gaps.

\subsection{False-Positive Audit: Methodology and Findings}
\label{sec:audit}
Given that cross-file inference is inherently more failure-prone than single-file pattern matching, we conducted an audit against ground truth, using a fixed random seed (20260703) and a repo-capped stratified sample ($\leq$2 findings per repository, spanning confidence tiers) to avoid single-repository clustering bias (Section~\ref{sec:audit-protocol}). Each sampled finding was independently traced against the actual source: for path claims, we verified whether the claimed reference (import, config key, named link) genuinely existed or was inferred from incidental co-location; for security findings, we verified the flagged content against its true context (executable statement vs.\ comment, docstring, or documentation).

\begin{table}[t]\centering\caption{False-positive audit across the three headline fix rounds (stratified $n{=}20$, seed 20260703, verified against source). A fourth round is discussed in Section~\ref{sec:eval-limits}.}\label{tab:audit}
\begin{tabular}{lrrl}\toprule
Engine state & FP rate & TP/20 & Root cause targeted \\\midrule
Initial (co-location edges) & 85\% & 3 & --- (baseline) \\
Post co-location fix & 90\%$^{\dagger}$ & 2 & Capability-word co-location edges \\
Post discovery-layer fix & 70\% & 6 & Docs/example/test as executable \\
Post precision fix & 60\% & 8 & Sink-label / docstring precision \\
\bottomrule\end{tabular}
\\[2pt]\footnotesize $^{\dagger}$The apparent rise from 85\% to 90\% reflects a change in what dominated the sample, not a worse engine: co-location edges (the initial dominant failure) were eliminated (0/18 failures in this round were co-location-based), revealing a previously-masked, larger class (documentation/test content misclassified as production) that the co-location fix did not target.
\end{table}


Each round targeted a single, empirically-identified dominant root cause, was verified with regression fixtures preventing recurrence, and was re-measured on the same seed for direct comparability. True positives nearly quadrupled (2 to 8) across the three headline rounds while total findings volume dropped 64\% (2{,}002 to 723) and cross-file paths dropped from an inflated 2{,}923 (pre-audit) to a reference-verified 520, indicating the reduction in volume reflects removal of false signal rather than loss of sensitivity.

\subsubsection{Remaining false-positive taxonomy}
At the final audited state (60\% FP, $n{=}20$), no remaining false positive belongs to either of the two most-recently-fixed root causes. We categorize the residual 12 false positives into six distinct, smaller-magnitude causes:
\begin{itemize}
\item \textbf{Docstring-as-prompt pipeline gap (2 cases):} our comment-stripping fix operates correctly at the rule level (verified via direct unit test), but the artifact parser extracts triple-quoted Python docstrings as standalone prompt strings \emph{before} rule evaluation, with delimiters removed --- making the extracted text indistinguishable from a genuine system prompt without full AST-level source-position tracking. This is a parser-level limitation, identified as future work.
\item \textbf{Keyword/store misattribution (2 cases):} capability inference still over-matches on some identifier patterns not covered by the credential-shape tightening applied in the precision fix.
\item \textbf{Sink label derived from remediation text rather than source content (2 cases):} in a small number of cases the reported sink reflects auto-generated remediation guidance rather than the underlying evidence, an internal labeling inconsistency.
\item \textbf{Single-cause classes:} CI-workflow-as-sink (1), i18n/config-file misclassification (1), unicode-rule precision (1), skill-instruction-as-sink (1), residual documentation outside recognized doc-directory conventions (1), and cross-file co-location residual (2) --- each independently verified and root-caused, none overlapping with the fixed classes.
\end{itemize}
We report this taxonomy in full rather than a single aggregate number because, in our experience building and auditing this tool, an unqualified false-positive percentage obscures more than it reveals: two engines with an identical 60\% FP rate can differ enormously in how tractable their remaining errors are. Ours are now concentrated in a long tail of independently fixable, well-scoped causes rather than one dominant, structurally difficult class --- evidence, we argue, that the audit-driven fixing methodology itself generalizes, even though full convergence to a low false-positive rate was not reached within our evaluation window.

\subsection{Scan Performance}
\label{sec:perf}
Table~\ref{tab:perf} reports scan performance. All scans are deterministic and single-threaded with zero LLM calls; identical repository states reproduce byte-identical output modulo timestamp and run-identifier fields (verified across 3 consecutive runs, Section~\ref{sec:determinism}).

\begin{table}[t]\centering\caption{Scan performance (seconds). Deterministic, single-threaded, no LLM calls.}\label{tab:perf}
\begin{tabular}{lrrrr}\toprule Corpus & Cases & Mean & Median & Max \\\midrule
Controlled & 22 & 0.4 & 0.2 & 1.6 \\
Wild & 37 & 10.9 & 6.6 & 38.7 \\
\bottomrule\end{tabular}\end{table}


\section{Limitations}
\label{sec:limitations}
We report limitations in three groups: detection scope gaps that affect what \sysname{} can observe at all, precision limitations quantified directly by our audit methodology (Section~\ref{sec:audit}), and evaluation limitations affecting how our results should be read.

\subsection{Detection scope}
\label{sec:scope-limits}
\textbf{Cross-file approval-policy ingestion is not implemented.} \sysname{}'s contextual verdict engine correctly adjusts severity based on approval signals found within a single artifact (e.g., an MCP server's \texttt{auto\_approve} field, exercised in the controlled corpus, Section~\ref{sec:controlled-eval}), but does not currently ingest a sibling approval-policy document and apply it to a separate skill or tool definition. A privileged capability documented as requiring human approval in an adjacent policy file is, at present, evaluated as if no such policy exists. This is a scope gap in artifact classification and control-state propagation, not a failure of the contextual verdict mechanism itself, which we verified functions correctly for in-artifact signals.

\textbf{Closure is downstream-only.} Repository discovery expands outward by following references from a seed set of recognized AI artifacts (prompts, skills, MCP configs) to the files they reference. It does not perform the reverse: acquiring code that references an AI artifact but is not itself referenced by one. In practice this means an entry-point file (e.g., an orchestrator that imports a skill definition) may be inventoried but excluded from analysis if nothing in the repository points back to it, truncating execution paths at the artifact rather than the true call-site origin.

\textbf{Docstring/prompt-string disambiguation is source-position-dependent, not semantic.} Our artifact parser extracts triple-quoted string literals matching prompt-like keyword heuristics as candidate prompt artifacts. We added AST-level detection to exclude genuine docstrings (a string that is the first statement of a module, class, or function body) from this extraction at the tree-sitter parsing stage. This fix is verified correct for module-level docstrings but incomplete for a secondary extraction path that operates on whole-file content independently of the tree-sitter walk; a small number of method-level docstrings and content-level keyword matches within docstrings can still surface as findings. Resolving this fully requires unifying the two extraction paths onto a single AST-aware pipeline, which we identify as future work.

\textbf{Format-gated secret and encoded-payload detection.} Hardcoded-secret detection requires a realistic provider-format credential shape (e.g., a recognized API key prefix and minimum length); encoded-payload detection requires a minimum length and a decode-to-injection-pattern condition. These gates are a deliberate precision/recall tradeoff --- verified during development to correctly ignore obviously-fake test fixtures --- but mean that short or atypically-formatted real secrets, and short or non-standard encoded payloads, are not flagged.

\subsection{Precision, quantified by audit}
\label{sec:precision}
Section~\ref{sec:audit} reports a three-round, audit-driven false-positive reduction from an initial 85\% ($n{=}20$, stratified sample, verified against source) to a final measured rate of 60\%, with true positives increasing from 3 to 8 across the same rounds. We report this rate directly rather than omitting it or reporting only aggregate finding counts, because we believe an unaudited prevalence claim from a static analyzer in this space is not trustworthy without independent verification, and our own early internal results (an unaudited claim that 89\% of sampled repositories contained a reachable path to shell execution) illustrate why: that figure, computed before our audit methodology was applied, was substantially inflated by exactly the false-positive classes the audit subsequently identified and partially corrected. \textbf{We do not report repository-level prevalence percentages in this paper's headline results; we report only audit-verified findings and the audit methodology itself as a contribution.}

The residual 60\% false-positive rate (Section~\ref{sec:audit}) decomposes into six identified causes, none dominant, each independently root-caused: a parser-level docstring-extraction gap (Section~\ref{sec:scope-limits}), sink-label misattribution from file content outside the classes already fixed, sink labels occasionally derived from auto-generated remediation text rather than source evidence, and several smaller single-cause classes (CI-workflow-as-sink, i18n/config misclassification, unicode-rule precision, residual documentation outside recognized directory conventions). We consider this taxonomy --- a long tail of small, independently-fixable causes rather than one dominant structural failure --- meaningfully different from, and more tractable than, the initial 85\% rate's single dominant cause (capability-word co-location), even though the numeric rate at evaluation time (60\%) remains too high for unsupervised production alerting without human review of Probable- and Potential-tier findings.

\subsection{Evaluation limitations}
\label{sec:eval-limits}
\textbf{Sample size.} Our false-positive audit uses $n{=}20$ per round, stratified by confidence tier with a per-repository cap to avoid clustering. At this sample size, tier-level false-positive rates (e.g., Confirmed-tier at 38\% in our final headline round) carry substantial uncertainty, particularly for the Confirmed tier where the underlying population is small. A subsequent, fourth fix round (targeting sink-label attribution from remediation text and the docstring-extraction gap) was audited on freshly re-cloned repository snapshots --- necessitated by an environment reset during development --- and reduced finding volume further (e.g., 723 to 620 findings on identical re-cloned snapshots) but produced a false-positive estimate of 65\% ($n{=}20$, 7 true positives) rather than a further reduction. We attribute the higher estimate to sampling variance at this sample size and to genuine content drift in actively-maintained repositories between clone snapshots, rather than to a regression in the underlying engine, which we verified by re-running both the pre- and post-fix builds against identical repository snapshots. We report this fourth round transparently as a distinct fix iteration whose targeted changes did not move the sampled rate, not as a re-verification of the 60\% state.

\textbf{Controlled-corpus construct validity.} Our controlled benchmark (Table~\ref{tab:controlled}, $P{=}R{=}F1{=}1.00$) was constructed by the authors, who also built the system under test. We report this result because it verifies mechanical correctness (path integrity, evidence-field completeness, determinism) against known ground truth, not as evidence of real-world precision --- the wild-corpus audit (Section~\ref{sec:audit}) is the paper's evidence for real-world precision, and it is deliberately independent of benchmark construction.

\textbf{Repository selection.} The 37-repository wild corpus was selected to span MCP server implementations, agent frameworks, and coding-assistant tooling among actively-maintained, popular open-source projects, but was not randomly sampled from the full population of public AI-agent repositories and may not be representative of that population, particularly of less mature or less widely-used projects.

\textbf{Baseline comparison scope.} Our baseline comparison (Section~\ref{sec:crossfile}) scopes both tools to each repository's AI-artifact surface rather than the full repository, after determining that whole-repository analysis caused the baseline to time out on large monorepos due to its dependency-vulnerability lookup stage, which is orthogonal to the skill-analysis capability under comparison. We believe this scoping produces a fairer comparison of the capability in question (per-artifact vs.\ cross-artifact reachability analysis) but note it is not an unmodified out-of-the-box comparison of either tool's default behavior on an arbitrary repository.

\section{Conclusion}
\label{sec:conclusion}
%% Not among the supplied drafts; brief placeholder to be authored.
\emph{[Conclusion to be authored: restate the cross-file capability contribution and the audit-as-methodology contribution, and the path from a 60\% audited FP rate toward production-usable precision.]}

\bibliographystyle{ACM-Reference-Format}
\bibliography{references}

%% ============================================================================
%% GENERATIVE-AI-USE DISCLOSURE (REQUIRED by AISec 2026; placed AFTER references,
%% does NOT count toward the page limit). HIGHEST-PRIORITY REMAINING GAP.
%% The author must write this paragraph directly (see summary doc, open decision
%% G1). Placeholder below is a starting scaffold ONLY and must be replaced.
%% ============================================================================
\section*{Disclosure of Generative-AI Use}
\emph{[PLACEHOLDER --- author to finalize.] The engineering and evaluation of this
work were carried out with substantial assistance from an AI coding assistant.
AI-assisted activities included implementation, debugging, construction of the
evaluation harness, and execution of the source-verification audit passes.
Human-directed decisions included the system architecture, root-cause
prioritization across fix rounds, scope decisions, and the framing of the
false-positive audit as a first-class contribution; all quantitative results
reported in this paper were verified by the authors against real repository
source. This paragraph must be reviewed and finalized by the authors to
accurately characterize the division of labor before submission.}

\end{document}
